\begin{proofbox}
\textbf{方法一｜纯代数：联立方程与判别式（严格充要证明）}\par
直线 $y=kx+m$ 代入椭圆 $b^2x^2+a^2y^2=a^2b^2$，整理为
\[
(b^2+a^2k^2)x^2+2a^2kmx+a^2(m^2-b^2)=0.
\]
由于首项系数 $b^2+a^2k^2>0$，它始终是真正的一元二次方程。其判别式为
\[
\begin{aligned}
\Delta
&=(2a^2km)^2-4(b^2+a^2k^2)a^2(m^2-b^2)\\
&=4a^2\Big[a^2k^2m^2-(b^2+a^2k^2)(m^2-b^2)\Big]\\
&=\boxed{4a^2b^2(a^2k^2+b^2-m^2)}.
\end{aligned}
\]
因为 $4a^2b^2>0$，故
\[
\begin{aligned}
\Delta=0&\iff m^2=a^2k^2+b^2 &&(\text{相切}),\\
\Delta>0&\iff m^2<a^2k^2+b^2 &&(\text{两交点}),\\
\Delta<0&\iff m^2>a^2k^2+b^2 &&(\text{无交点}).
\end{aligned}
\]
这同时证明了\textbf{公式的充分性与必要性}以及全部位置关系。

\medskip
\textbf{方法二｜纯几何：椭圆伸缩成单位圆（最容易想象）}\par
作可逆伸缩变换 $x=aX,\ y=bY$。椭圆变成单位圆
\[
X^2+Y^2=1,
\]
直线变成
\[
-akX+bY-m=0.
\]
可逆的一一对应变换不会改变交点的个数，所以原直线与椭圆相切，当且仅当新直线与单位圆相切。单位圆的圆心为 $O(0,0)$，半径为 $1$。由点到直线距离公式，
\[
d(O,l')=\frac{|m|}{\sqrt{(-ak)^2+b^2}}.
\]
因此
\[
\text{相切}\iff d(O,l')=1
\iff |m|=\sqrt{a^2k^2+b^2}
\iff \boxed{m^2=a^2k^2+b^2}.
\]
同理，距离小于 $1$ 为相交，大于 $1$ 为相离。此证明不需要解一元二次方程。

\medskip
\textbf{方法三｜数形结合：参数方程＋辅助角公式}\par
令椭圆上动点 $P(a\cos\theta,b\sin\theta)$，$0\le\theta<2\pi$。点 $P$ 在直线 $y=kx+m$ 上，当且仅当
\[
m=-ak\cos\theta+b\sin\theta.
\]
令 $R=\sqrt{a^2k^2+b^2}>0$。可取角 $\varphi$ 满足
\[
\cos\varphi=\frac{b}{R},\qquad
\sin\varphi=\frac{-ak}{R},
\]
于是
\[
m=R\sin(\theta+\varphi).
\]
因为 $\sin(\theta+\varphi)$ 的值域是 $[-1,1]$，所以：
\begin{itemize}
    \item $|m|<R$：一周内有两个不同的 $\theta$，对应两个交点；
    \item $|m|=R$：一周内仅有一个 $\theta$，对应唯一交点，即相切；
    \item $|m|>R$：不存在交点。
\end{itemize}
所以 $\boxed{m^2=R^2=a^2k^2+b^2}$。\textbf{直观意义}：相切就是 $y-kx$ 取到最大值或最小值。

\medskip
\textbf{方法四｜从切点反推：椭圆切线方程法}\par
设有限斜率切线在 $T(x_0,y_0)$ 处与椭圆相切，则 $y_0\ne0$，且
\[
\frac{x_0x}{a^2}+\frac{y_0y}{b^2}=1,
\qquad
\frac{x_0^2}{a^2}+\frac{y_0^2}{b^2}=1.
\]
把切线方程整理成斜截式，得到
\[
k=-\frac{b^2x_0}{a^2y_0},\qquad m=\frac{b^2}{y_0}.
\]
于是
\[
\begin{aligned}
m^2-a^2k^2
&=\frac{b^4}{y_0^2}\left(1-\frac{x_0^2}{a^2}\right)\\
&=\frac{b^4}{y_0^2}\cdot\frac{y_0^2}{b^2}
=\boxed{b^2}.
\end{aligned}
\]
即 $m^2=a^2k^2+b^2$。这先证明必要性；反过来，若该等式成立，设
\[
T\left(-\frac{a^2k}{m},\frac{b^2}{m}\right),
\]
则 $m\ne0$，而且
\[
\frac{x_T^2}{a^2}+\frac{y_T^2}{b^2}
=\frac{a^2k^2+b^2}{m^2}=1.
\]
在 $T$ 处的切线
\[
\frac{x_Tx}{a^2}+\frac{y_Ty}{b^2}=1
\]
化为 $-kx+y=m$，正是原直线。故充分性也成立，并且顺便求出了\textbf{切点坐标}。

\medskip
\textbf{方法五｜最短推导：向量投影／柯西不等式}\par
椭圆上的点满足
\[
\left(\frac xa\right)^2+\left(\frac yb\right)^2=1.
\]
由柯西不等式，
\[
\begin{aligned}
(y-kx)^2
&=\left[(-ak)\frac xa+b\frac yb\right]^2\\
&\le\big[(ak)^2+b^2\big]
\left[\left(\frac xa\right)^2+\left(\frac yb\right)^2\right]\\
&=a^2k^2+b^2.
\end{aligned}
\]
等号当且仅当单位向量 $(x/a,y/b)$ 与向量 $(-ak,b)$ 同向或反向。此时 $y-kx$ 恰好取到两端极值 $\pm\sqrt{a^2k^2+b^2}$，每个极值对应唯一切点。
因此
\[
\boxed{\text{相切}\iff m^2=a^2k^2+b^2.}
\]
\textbf{记忆核心}：$y-kx$ 是向量 $(-ak,b)$ 在单位圆上的投影，最大绝对值就是这个向量的长度。

\medskip
\textbf{拓展｜一次解决所有直线形式}\par
对 $Ax+By+C=0$，$A,B$ 不同时为零，作同样伸缩 $x=aX,\ y=bY$，得到
\[
aAX+bBY+C=0.
\]
该直线与单位圆相切，当且仅当
\[
\frac{|C|}{\sqrt{a^2A^2+b^2B^2}}=1
\iff\boxed{C^2=a^2A^2+b^2B^2}.
\]
这个公式\textbf{不用区分斜率是否存在}。取 $A=-k,B=1,C=-m$ 就恢复斜率式；取 $B=0$ 就得到竖直切线 $x=\pm a$。
\end{proofbox}
